\documentclass[11pt]{article}

\usepackage{fullpage}
\usepackage{enumitem}
\usepackage{amsmath,amsthm,amsfonts,amssymb,mathtools,latexsym}
\usepackage{xcolor,graphicx,hyperref,float}
\usepackage{tikz}
\usetikzlibrary{arrows.meta}
\usepackage[ruled,vlined,linesnumbered]{algorithm2e}

\pdfpagewidth 8.5in
\pdfpageheight 11in
\oddsidemargin 0in
\evensidemargin 0in

\providecommand{\tightlist}{\setlength{\itemsep}{0pt}\setlength{\parskip}{0pt}}
\providecommand{\squigrightarrow}{\rightsquigarrow}
\providecommand{\set}[1]{\{#1\}}

\begin{document}
\begin{center}
\Large{\textbf{Homework 9}}
\end{center}

\hypertarget{recurrences}{%
\section{1. Recurrences}\label{recurrences}}

\hypertarget{a-a-terminates-when-n-is-a-non-negative-int}{%
\subsection{\texorpdfstring{(a) A terminates when \(n\) is a
non-negative
\texttt{int}}{(a) A terminates when n is a non-negative int}}\label{a-a-terminates-when-n-is-a-non-negative-int}}

Proof by Strong Induction

\begin{itemize}
\tightlist
\item
  Base case: \(n=0\)

  \begin{itemize}
  \tightlist
  \item
    the function terminates immediately if \(n=0\) that return \(0\).
  \end{itemize}
\item
  Inductive step

  \begin{itemize}
  \tightlist
  \item
    every odd integer input \(n>1\) returns a recursive call to even
    input \(2(n-1)\).
  \item
    after dividing by 2 in even integer input case, \(n-1\).
  \item
    So, each odd/even call reduces the non-negative integer input by
    \(1\).
  \item
    after \(2 \times n\) where \(n\) iterations where \(n\) is the
    initial non-negative integer input, the function will eventually
    return \(n=0\).
  \end{itemize}
\end{itemize}

Therefore, \(A\) terminates when its input \(n\) is a non-negative
integer.

\hypertarget{b-recurrence-for-the-running-time-of-a}{%
\subsection{(b) recurrence for the running time of
A}\label{b-recurrence-for-the-running-time-of-a}}

\(T(n) = T(n/2) + \Theta(1)\): even case
\(T(n) = T(2(n-1)) + \Theta(1)\): odd case

\hypertarget{c}{%
\subsection{(c)}\label{c}}

\[\begin{array}{ll}
T(n/2) + \Theta(1)\\
a=1, b=2, d=0\\
0 = log_2(1) \rightarrow T(n) = \Theta(n^0logn)\\
= \Theta(logn)\\
\\
T(2(n-1)) + \Theta(1)\\
= \Theta(1) + 2\Theta(1) + 4\Theta(1) +...+2^{n-1}\Theta(1)\\
= 2^n\Theta(1) = \Theta(2^n)
\end{array}
\]

\hypertarget{d}{%
\subsection{(d)}\label{d}}

\begin{enumerate}
\def\labelenumi{\arabic{enumi}.}
\tightlist
\item
  depth: \(O(logn)\)
\item
  the number of leaves: \(O(1)\); exactly 1 leaf left of base case
\end{enumerate}

\hypertarget{e}{%
\subsection{(e)}\label{e}}

It computes the number of \(1\)s of the input \(n\) in binary
expression.

\begin{itemize}
\tightlist
\item
  e.g.~\(10\) = \texttt{0b1010}

  \begin{itemize}
  \tightlist
  \item
    \(A(10) \text{ where n = 10}: 2\)
  \item
    \(A(10) \rightarrow A(5) \rightarrow A(8) + 1(\text{ this is }\Theta(1) \text{operation}) \rightarrow A(4) \rightarrow A(2) \rightarrow A(1) + 1 \rightarrow A(0) (\text{base case reached})\)
  \item
    count the number of \(1\)s: \(1+1=2\).
  \end{itemize}
\end{itemize}

\hypertarget{carpets}{%
\section{2. Carpets}\label{carpets}}

\hypertarget{algorithm}{%
\subsection{1. Algorithm}\label{algorithm}}

\hypertarget{plain-english-version}{%
\subsubsection{plain English version}\label{plain-english-version}}

\hypertarget{io}{%
\paragraph{I/O}\label{io}}

\begin{itemize}
\item
  Input: \(n=2^k\) (size of the room) and coordinate
  \((x,y) \in \set{1,...,n} \times \set{1,...,n}\) of the \(1 \times 1\)
  covered square.
\item
  Output: list of carpet pieces to fill up \(2^k \times 2^k\) floor
\end{itemize}

\hypertarget{divide-and-conquer-recursive-call}{%
\paragraph{Divide and conquer recursive
call}\label{divide-and-conquer-recursive-call}}

\begin{enumerate}
\def\labelenumi{\arabic{enumi}.}
\tightlist
\item
  divide the current square floor in 4 quadrants
\item
  we decide which quadrant we will recursively call by getting which
  quadrant the input \(1 \times 1\) covered resides
\item
  we cover the floor of the other three quadrant without the input
  covered tile with the matching carpet piece.
\item
  as we cover the floor with the matching carpet pieces depending on the
  location of given \(1 \times 1\) cover resides, we update the
  coordinate and its matching pieces in the list
\item
  we process the whole quadrants recursively with the initially added 1
  cover for the other 3 quadrants with L-shaped cover piece at the
  middle. stop edge case for the recursion would be to reach
  \(1 \times 1\) grid.
\item
  return the cover piece list.
\end{enumerate}

\hypertarget{pseudocode}{%
\subsubsection{Pseudocode}\label{pseudocode}}

\begin{verbatim}
carpet(n: int, cx: int, cy: int):
    /*IO
    Input: n: size of the floor; power of 2
    cx: x coordianatce of the already-covered floor
    cy: y coordinate of the already-covered floor

    Output: list of tuples that contains: (x, y, L-shaped carpet piece)
    */

    // 1. base case: we stop when we don't have space to put L-shaped carpet piece
    if n = 1: // size of 1
        return empty list

    center = n / 2 // half the size so make 4 quadrants
    quadlist = {
        TL = (center, center)
        TR = (center + 1, center)
        BL = (center, center + 1)
        BR = (center + 1, center + 1)
    }
    // find which quad does (cx,cy) reside
    if cx <= center:
        if cy <= center:
            quad = TL
        else:
            quad = BL
    else:
        if cy <= center:
            quad = TR
        else:
            quad = BR
    // locating other three quads; need to process all four quadrants
    opposite = {TL: BR, TR: BL, BL: TR, BR: TL} // already-covered: opposite direction to fillup
    fillup = quadlist[opposite[quad]] // get the coordinates of the directional l-shape cover piece

    L_tiles = [(fillup.cx, fillup.cy, quad + "direction")] // output L tile

    subgrid = {
        TL: (0,0,quadlist[TL]),
        TR: (center, 0, quadlist[TR]),
        BL: (0, center, quadlist[BL]),
        BR: (center,center, quadlist[BR])
    }
   // recursion order: first, quad that the initially given $1 \time 1$ cover resides
   // then, do fillups at quadrants clockwise-neighbor
    recurse_order = [quad] + [for i in subgrid: subgrid[i].key if subgrid[i].key == quad: continue]
    
    // divide subproblems by quad
    for q' in recurse_order:
        (x_offset, y_offset, center_cell) = subgrid[q']

        // recurse until we reach the base case
        if q' = quad:
            sub_x = cx - x_offset
            sub_y = cy - y_offset
        else:
            sub_x = center_cell.cx - x_offset
            sub_y = center_cell.cy - y_offset

        // recurse on smaller subproblems
        sub_tiles = carpet(center, sub_x, sub_y)

        // append the l shape tiles in the output
        for (x, y, dir) in sub_tiles:
            L_tiles.append(x + x_offset, y + y_offset, dir)
    return L_tiles
\end{verbatim}

\hypertarget{proof}{%
\subsection{proof}\label{proof}}

To prove the algorithm \texttt{carpet(n,\ cx,\ cy)} is correct:

\begin{itemize}
\tightlist
\item
  for each recursive call:

  \begin{itemize}
  \tightlist
  \item
    we get a square subgrid(subfloor) with exactly one already-covered
    cell
  \item
    return tiles that cover the other cells exactly once
  \end{itemize}

  Given \(n=2^k\):

  \begin{enumerate}
  \def\labelenumi{\arabic{enumi}.}
  \tightlist
  \item
    base case: \(1 \times 1\) subgrid has one already-covered cell that
    fills up the whole subgrid: return no tiles
  \item
    Inductive step: assume this holds as well for \(n/2 \times n/2\)
    subgrids.

    \begin{itemize}
    \tightlist
    \item
      because the algorithm checks four \(n/2 \times n/2\) subgrids to
      check which subgrid has the input already-covered cell and put
      exactly one \(1 \times 1\) cover for every other subgrid, the
      claim for each recursive call holds.
    \item
      for each \(n/{2^k} \times n/{2^k}\) subgrids that has exactly one
      \(1 \times 1\) cover, the leftover grid to fillup, considering we
      only have L-shaped covers, is
      \(\frac{(n/{2^k} \times n/{2^k})-1}{3}\). This is exactly the same
      with the number of carpet pieces that friend gave us. Therefore,
      there is no overlap for filling up the whole four subgrid.
    \end{itemize}
  \end{enumerate}

  Therefore, the algorithm is always correct
\end{itemize}

\hypertarget{recurrence-for-the-worst-running-time-in-terms-of-n-and-the-running-time}{%
\subsection{\texorpdfstring{recurrence for the worst running time in
terms of \(n\) and the running
time}{recurrence for the worst running time in terms of n and the running time}}\label{recurrence-for-the-worst-running-time-in-terms-of-n-and-the-running-time}}

\(T(n) = 4T(n/2) + \Theta(n^2): a = 4, b = 2, d = 2\)

\begin{itemize}
\item
  we divide into 4 subproblems (quadrants)
\item
  we make the problem size by half
\item
  tile list conversion: \(n^2\) (nested for loop)
\end{itemize}

Solve by Master method
\(d=2=log_2(4) \rightarrow T(n) = \Theta(n^2logn)\)
\end{document}
